\documentclass[11pt]{article}
\usepackage[margin=1in]{geometry}
\usepackage{amsmath,amssymb,booktabs,enumitem,microtype}
\usepackage[colorlinks=true,linkcolor=blue,urlcolor=blue,citecolor=blue]{hyperref}

\setlist{leftmargin=*,itemsep=3pt}
\newcommand{\major}{\textbf{Major}}
\newcommand{\moderate}{\textbf{Moderate}}
\newcommand{\minor}{\textbf{Minor}}

\title{Adversarial Mathematical Assessment of\
\texttt{BEMOCOneEnergy.tex}}
\author{}
\date{13 July 2026}

\begin{document}
\maketitle

\begin{abstract}
This report independently audits the manuscript \texttt{BEMOCOneEnergy.tex},
which claims
\[
 \frac43N^2-\sum_{x\ne y\in\mathcal P_N}|x-y|=\Theta(\sqrt N)
\]
for the Beltr\'an--Etayo--Marzo--Ortega-Cerd\`a point set.  The exact
decomposition, the within-parallel calculation, the uniform Fourier estimate,
the greatest-common-divisor summation, and the Stolarsky--Beck lower bound are
essentially correct.  The proof nevertheless does not yet pass referee-level
scrutiny: its decisive weakly singular latitude-block lemma is only sketched,
and one of the geometric comparisons used to claim global uniformity is false
for comparable-scale bands in opposite hemispheres.  This appears repairable
and is not a counterexample to the theorem, but until the missing antipodal
case and the asserted derivative bounds are proved, the upper bound
$D_N=O(\sqrt N)$ remains incomplete.
\end{abstract}

\section*{Post-audit revision note}

After this report was written, \texttt{BEMOCOneEnergy.tex} was revised to
implement the repair proposed below.  The new kernel lemma separates
same-hemisphere near-diagonal pairs, near-equatorial cross-hemisphere pairs,
genuinely antipodal pairs, unequal latitude scales, and the exceptional
central band that may straddle the equator.  The exceptional equatorial gcd
sum and the minor citation and notation issues were also corrected.  An
independent re-audit of the revised manuscript returned \textbf{PASS}.  The
report below is retained as a record of the defects in the earlier version
and of how the repair was derived.

\section{Summary verdict}

Write
\[
 \Sigma_N=\sum_{x\ne y\in\mathcal P_N}|x-y|,
 \qquad D_N=\frac43N^2-\Sigma_N.
\]
The manuscript does not currently prove $D_N=\Theta(\sqrt N)$.  The lower
bound is valid, and all parts of the proposed upper bound except the latitude
quadrature estimate withstand the audit.  The remaining issue is a
\major{} gap in the proof as written, rather than evidence that the theorem is
false.

\begin{center}
\begin{tabular}{p{0.30\textwidth}p{0.17\textwidth}p{0.43\textwidth}}
\toprule
Component & Verdict & Reason\\
\midrule
Construction and cardinality & Correct, minor citation error & The formulas
agree with the BEMOC construction, but the admissible choice occurs in
Lemma~1.6, not Lemma~1.7.\\
Exact decomposition and signs & Correct & The Legendre coefficients and the
positive-definiteness argument give $A_N,D_N\ge0$.\\
Within-ring term & Correct & The cotangent identity and the square-subsequence
constant check.\\
Cross-ring term & Essentially correct & The bounded-variation Fourier bound
is uniform, and the resulting gcd sum has the claimed order; exceptional
equatorial rings need one explicit sentence.\\
Latitude block estimate & Major gap & Uniform derivative estimates are
asserted rather than established, and the separation comparison used for
opposite hemispheres is false.\\
Universal lower bound & Correct & Stolarsky's normalization and Beck's
$N^{-3/2}$ squared-discrepancy bound yield $D_N\gtrsim\sqrt N$.\\
\bottomrule
\end{tabular}
\end{center}

Accordingly, the conditional project of extending the argument to other
Riesz exponents should not yet be undertaken: it would inherit precisely the
unproved kernel estimate on which the present exponent depends.

\section{Components that survive the audit}

\subsection{Construction}

For $M=\lfloor\sqrt{N/4}\rfloor$ one has
\[
 r_M=N-4M^2+6M-2,
\]
so $r_M\asymp M$ throughout $4M^2\le N<4(M+1)^2$.  The ring populations in
the manuscript agree with the Simpson-type atomization in Section~4 of
BEMOC.  The text should distinguish the canonical phases selected in BEMOC
from the stronger assertion made here, namely uniformity under arbitrary
independent rotations of the polygons.

The opening citation says ``Lemma~1.7.''  In the published and arXiv versions,
the displayed choice of $M,r_1,\ldots,r_M$ is Lemma~1.6; item 1.7 is a remark.

\subsection{Exact decomposition and spectral sign}

For a ring of radius $\rho$, its continuous angular self-energy is
$4\rho w^2/\pi$.  Adding and subtracting all continuous ring--ring energies
therefore gives the stated exact identity
\[
 D_N=A_N+B_N+C_N.
\]
The absolutely summable Legendre expansion
\[
 \sqrt{2-2u}=\frac43-\sum_{n\ge1}
 \frac4{(2n-1)(2n+3)}P_n(u)
\]
and the addition theorem imply that the energy of every positive measure of
mass $N$ is at most $4N^2/3$.  Thus $A_N\ge0$ and $D_N\ge0$ as claimed.

\subsection{Within-ring contribution}

For $w$ equally spaced points on a parallel of radius $\rho$,
\[
 \Sigma^{\rm self}=2\rho w\cot\frac{\pi}{2w},
\]
and hence
\[
 \frac{4\rho w^2}{\pi}-\Sigma^{\rm self}
 =\frac{\pi\rho}{3}+O(\rho w^{-2}).
\]
The BEMOC relations $w_p\asymp\sqrt N\,\rho_p$ and
$\sum_p\rho_p=O(\sqrt N)$ give $0\le B_N\le C\sqrt N$.  Along $N=4m^2$,
the Riemann-sum calculation yields
\[
 \frac{B_N}{\sqrt N}\longrightarrow
 \frac{2\pi(2\sqrt2-1)}9.
\]
Line 147 contains only a typesetting error: ``$cot xle x^{-1}$'' should be
``$\cot x\le x^{-1}$''.

\subsection{Cross-ring aliasing}

For $G(\theta)=\sqrt{A-B\cos\theta}$ with $A\ge B\ge0$, the bound
$|G'|\le C\sqrt B$ is uniform, and $G'$ has a uniformly bounded number of
monotonicity intervals.  Including the cusp when $A=B$, this gives
\[
 \|DG'\|_{\rm TV}\le C\sqrt B,
 \qquad |\widehat G(n)|\le C\sqrt B\,n^{-2}.
\]
The least-common-multiple aliasing identity consequently gives, for ring
populations $q,r$,
\[
 |\mathcal E(q,r)|\le
 C\sqrt{\rho_q\rho_r}\frac{(q,r)^2}{qr}.
\]
Using $q\asymp M\rho_q$ and the bounded multiplicity of the actual BEMOC
populations reduces the total to
\[
 \frac CM\sum_{q,r\le CM}\frac{(q,r)^2}{\sqrt{qr}}.
\]
Jordan's identity and $J_2(a)\le a^2$ give an $O(M^2)$ gcd sum and hence
$|C_N|=O(M)=O(\sqrt N)$.  The $O(1)$ rings involving $r_M$ should be written
out separately, but (4.1) gives the same admissible total bound.

\subsection{Stolarsky--Beck lower bound}

With normalized surface area and ordinary Lebesgue measure in the cap-height
variable, the factor $4$ in the displayed Stolarsky identity is correct.
Beck's lower bound for squared spherical-cap $L^2$ discrepancy has order
$N^{-3/2}$; multiplication by $N^2$ gives
\[
 D_N\ge c\sqrt N.
\]

\section{The unresolved latitude estimate}

Let $\mu_j$ be the signed error of the Simpson/midpoint atomization on band
$B_j$.  The identities
\[
 \int d\mu_j=\int t\,d\mu_j(t)=0,
 \qquad \|\mu_j\|_{\rm TV}\lesssim r_j,
 \qquad |B_j|\asymp r_j/N
\]
are correct and are exactly the cancellation needed for a viable proof.
Likewise, the local diagonal singularity of the reduced distance kernel has
the form
\[
 \frac{(s-t)^2}{2\pi R^3}\log\frac{R^2}{|s-t|}
\]
up to smoother terms when $R\asymp\sqrt{1-s^2}\asymp\sqrt{1-t^2}$.

The difficulty is that lines 334--399 do not establish the global block
lemma at the stated level of uniformity.

\subsection{A false separation comparison}

Lines 375--384 use
\[
 |s-t|\asymp \ell a_j,
 \qquad \ell=|j-k|,
 \qquad a_j\asymp d_j/N,
\]
for all comparable-scale bands.  Take the north band $j=2$ and its south
mirror $k=2M-2$.  Then
\[
 d_j=d_k=2,qquad \ell=2M-4,qquad a_j\asymp N^{-1}.
\]
Thus $\ell a_j\asymp N^{-1/2}$, whereas points in the two bands satisfy
\[
 |s-t|=2-O(N^{-1}).
\]
The asserted comparability is therefore false.  The later sentence about
``cross-equatorial neighbors'' handles only bands close to the equator; it
does not cover comparable-scale bands close to opposite poles.

This does not disprove the block estimate.  In fact the kernel is smoother
in the genuinely antipodal regime, so a stronger estimate is expected.  It
does mean that the proof printed in the manuscript does not cover all pairs
summed in the necklace bound.

\subsection{Derivative bounds are stated, not proved}

Equation (5.8) introduces a term $L(s,t)$ described as ``affine in each
variable modulo a remainder,'' while the required bounds on that remainder
and on four dimensionless derivatives are asserted after a reference to the
chain rule.  This is the analytic heart of the proof.  To apply the
two-variable Peano formula uniformly, the manuscript must state and prove a
precise derivative lemma in the following distinct regimes:
\begin{enumerate}
\item comparable bands in one hemisphere near the diagonal;
\item comparable bands in opposite hemispheres but close to the equator;
\item genuinely antipodal pairs;
\item unequal latitude scales, including one polar block and one bulk block.
\end{enumerate}
The rescaling $s=1-x/N$, $t=1-y/N$ only treats pairs near the same pole.  By
itself it does not justify interactions between a polar block and all other
latitudes.

\section{A viable repair}

The existing near-diagonal calculation can be retained.  Split the omitted
opposite-hemisphere pairs into two regions.

Near the equator, $d_j,d_k\asymp M$ and the vertical separation is comparable
to $|j-k|/M$, so the current Peano estimate can be made rigorous after fixing
the derivative lemma.  Away from the equator, the two heights stay a
controlled distance from the diagonal singularity.  For the polar mirror
case, set
\[
 s=1-x/N,\qquad t=-1+y/N.
\]
The averaged kernel has an ordinary expansion
\[
 F(s,t)=2-\frac{x+y}{2N}+N^{-2}P_2(x,y)+\cdots.
\]
Zero mass and zero first moment annihilate all terms that do not have degree
at least two in each block variable.  A four-derivative Peano estimate then
gives a bound stronger than (5.3).  The same argument, with compactness away
from the diagonal, covers the remaining antipodal region.

Once these cases are written as a uniform lemma, summing (5.3)--(5.4) does
indeed give $A_N=O(\sqrt N)$.  Until then, that bound and consequently the
upper half of the main theorem remain unproved.

\section{Status of a general-Riesz extension}

The exponent $s=-1$ is unusually forgiving: after angular averaging its
latitude kernel has only the weak singularity
$(u-v)^2\log|u-v|$, and the polygonal aliasing has $n^{-2}$ Fourier decay.
For a general Riesz kernel $|x-y|^{-s}$ both singular orders change.  Any
extension would therefore require a parameter-dependent version of the very
latitude lemma that is currently incomplete, plus a new uniform aliasing
estimate.  Formally differentiating or substituting an exponent into the
present proof would repeat the defect identified above.

Because the requested extension was conditional on the manuscript passing
adversarial review, this report does not claim such an extension.  The sound
next step is to repair and independently verify the four-regime latitude
lemma first.

\begin{thebibliography}{9}
\bibitem{BEMOC}
C.~Beltr\'an, U.~Etayo, J.~Marzo and J.~Ortega-Cerd\`a,
\emph{A sequence of polynomials with optimal condition number},
J. Amer. Math. Soc. \textbf{34} (2021), 219--244;
\href{https://doi.org/10.1090/jams/956}{doi:10.1090/jams/956},
\href{https://arxiv.org/abs/1903.01356}{arXiv:1903.01356}.

\bibitem{Beck}
J.~Beck, \emph{Sums of distances between points on a sphere---an application
of the theory of irregularities of distribution to discrete geometry},
Mathematika \textbf{31} (1984), 33--41;
\href{https://doi.org/10.1112/S0025579300010639}{doi:10.1112/S0025579300010639}.

\bibitem{Stolarsky}
K.~B.~Stolarsky, \emph{Sums of distances between points on a sphere. II},
Proc. Amer. Math. Soc. \textbf{41} (1973), 575--582;
\href{https://doi.org/10.1090/S0002-9939-1973-0333995-0}{doi:10.1090/S0002-9939-1973-0333995-0}.
\end{thebibliography}

\end{document}
